\bookchapter{第一章：你好，我的私人電腦}

\chapterintro{你的手機，其實是你口袋裡最私人的超級計算機。它比當年把人類送上月球的計算機強大無數倍，卻安靜地躺在你的掌心。今天，咱們就來好好認識這位“沉默的巨人”。}

\booksection{一、手機硬體“五臟六腑”：深度解析}

手機雖小，五臟俱全。讓我們像解剖課一樣，逐個認識這些精密的元件。

\booksubsection{SoC（系統級晶片）：手機的“大腦+心臟”}

SoC全稱System on Chip，翻譯過來就是“把整個系統塞進一塊晶片”。它不只是CPU（中央處理器），而是把CPU、GPU（圖形處理器）、NPU（神經網路處理器）、ISP（影象訊號處理器）、基帶（負責通訊）等等全部整合在一起。

這就像把一個城市的政府大樓、發電廠、通訊中心、警察局全部建在一棟樓裡，效率極高。

\booksubsection{CPU：通用計算的核心}

CPU是SoC中最核心的元件，負責處理絕大多數計算任務。手機CPU採用ARM架構，與電腦的x86架構不同，ARM架構以精簡指令集（RISC）為特色，功耗更低，更適合移動裝置。

\textbf{CPU大小核設計}：現代手機CPU普遍採用“大小核”架構（ARM的big.LITTLE技術），將核心分為不同效能等級：

\begin{itemize}
  \item \textbf{超大核/大核}：效能最強，負責重度任務。如遊戲載入、應用啟動、複雜計算。頻率最高（3GHz+），但功耗也最大。
  \item \textbf{中核}：效能與功耗平衡，處理日常中等負載任務。刷微博、看影片、多工切換主要靠它。
  \item \textbf{小核/能效核}：功耗最低，處理輕量級後臺任務。待機、聽音樂、接收訊息等，幾乎不費電。
\end{itemize}

三個等級的核心協同工作：系統排程器根據任務負載自動分配——輕量任務交給小核省電，重活交給大核猛幹，幹完再切回小核。這就是為什麼手機既能流暢打遊戲，又能待機一整天。

\textbf{各品牌架構特點}：

\begin{itemize}
  \item \textbf{高通驍龍8 Elite}：採用Oryon自研CPU架構（源自Nuvia收購），2個超級核（4.32GHz）+ 6個效能核（3.53GHz），全部基於臺積電第二代3nm製程。徹底告別ARM公版架構，單核和多核效能均實現跨越式提升。GPU為Adreno 830，採用切片架構，圖形效能提升40\%，Hexagon NPU算力約45 TOPS，AI推理效能提升45\%。（注：三星Galaxy定製版頻率更高，超級核可達4.47GHz。）
  \item \textbf{聯發科天璣9400}：1個超大核（Cortex-X925，3.62GHz）+ 3個效能核（Cortex-X4，3.3GHz）+ 4個能效核（Cortex-A720，2.4GHz），全大核架構延續，臺積電第二代3nm製程。GPU為Immortalis-G925（12核），峰值效能提升41\%，光線追蹤效能提升40\%。AI處理器為第八代NPU 890，支援端側大模型執行、端側LoRA訓練和端側影片生成。
  \item \textbf{蘋果A18 Pro}：2個效能核（約4.04GHz）+ 4個能效核（約2.4GHz），臺積電第二代3nm製程。6核GPU支援硬體級光線追蹤，16核NPU（約35 TOPS）為端側AI和Apple Intelligence提供強大算力。蘋果自研架構在單核效能和能效比上依然領先，統一記憶體架構為AI推理提供高頻寬優勢。
  \item \textbf{華為麒麟9020}：搭載於Mate 70系列，採用自研泰山CPU架構，NPU和ISP能力持續增強，支援5.5G和衛星通訊，代表國產晶片的最新突破。
\end{itemize}

\booksubsection{GPU：圖形渲染引擎}

GPU（Graphics Processing Unit）專門負責圖形和平行計算任務。它的核心數量遠多於CPU（動輒數百上千個核心），但每個核心的功能相對簡單，適合同時處理大量相似的計算。

\textbf{圖形渲染流程}：

\begin{enumerate}
  \item \textbf{頂點處理}：計算3D模型中每個頂點的位置座標。
  \item \textbf{光柵化}：將3D圖形轉換為2D螢幕上的畫素點。
  \item \textbf{畫素著色}：為每個畫素計算顏色、紋理、光照效果。
  \item \textbf{輸出合成}：將所有圖層合成最終畫面，送到螢幕顯示。
\end{enumerate}

\textbf{GPU對日常使用的影響}：不只是遊戲！介面滑動、動畫效果、影片播放、照片濾鏡、AR應用都依賴GPU。GPU效能弱的手機，滑動介面可能掉幀，遊戲更是直接卡頓。

\textbf{各品牌GPU}：高通使用自研Adreno GPU；聯發科使用ARM Mali GPU；蘋果使用自研GPU，效能和能效都是行業標杆。

\booksubsection{NPU：AI加速引擎}

NPU（Neural Processing Unit）是專門為神經網路運算設計的處理器。相比CPU和GPU，NPU在處理AI任務時效率更高、功耗更低。

\textbf{神經網路運算方式}：NPU的核心運算是矩陣乘法和卷積運算——這正是神經網路中最基礎、最頻繁的計算。NPU透過硬體級最佳化，可以並行處理大量矩陣運算，速度比CPU快幾十倍，功耗卻低得多。

\textbf{AI處理流程}：

\begin{enumerate}
  \item \textbf{資料輸入}：攝像頭影象、麥克風音訊、感測器資料等。
  \item \textbf{預處理}：調整資料格式和尺寸，適配神經網路輸入要求。
  \item \textbf{推理計算}：NPU載入訓練好的模型，對輸入資料進行前向計算，輸出結果。
  \item \textbf{後處理}：將AI輸出轉換為可用的結果，如人臉座標、文字內容、場景標籤。
\end{enumerate}

\textbf{實際應用場景}：

\begin{itemize}
  \item \textbf{計算攝影}：夜景合成、人像虛化、HDR處理、AI美顏。
  \item \textbf{語音助手}：語音識別、自然語言理解、語音合成。
  \item \textbf{實時翻譯}：AI字幕、對話翻譯。
  \item \textbf{智慧識別}：文字識別（OCR）、物體識別、場景識別。
  \item \textbf{端側大模型}：最新旗艦晶片已支援在手機端執行小型AI大模型（7B引數級別），實現離線AI對話、文件摘要、智慧寫作等功能。驍龍8 Elite的Hexagon NPU算力約45 TOPS，天璣9400的NPU 890同樣支援端側大模型（LLM推理速度提升80\%），蘋果A18 Pro的16核NPU提供約35 TOPS算力。
  \item \textbf{AI影象處理}：AI消除（一鍵去除照片中的路人）、AI擴圖（擴充圖片邊界）、AI風格轉換、文字生成圖片等。
\end{itemize}

\booksubsection{ISP：影像訊號處理器}

ISP處理攝像頭傳來的原始資料，把光線訊號變成好看的照片。

\textbf{效能怎麼看？}看型號！蘋果A18 Pro、高通驍龍8 Elite、聯發科天璣9400、華為麒麟9020都是旗艦級別。中端看驍龍8s Gen3/7+ Gen3、天璣8400。別隻看核心數，架構和製程（奈米數）更重要。3nm製程是目前最先進的，能效比遠優於5nm和4nm。

\booksubsection{運存（RAM）：臨時工作臺}

運存就像你的辦公桌，越大，能同時攤開的檔案就越多。12GB是現在的主流，16GB適合重度使用者，24GB以上則是遊戲玩家和極客的選擇。

運存型別也很重要：LPDDR5比LPDDR4X更快更省電。蘋果雖然運存不大，但系統最佳化好，所以看起來“夠用”。

\textbf{運存不夠會怎樣？}App會被系統殺掉後臺，切換時需要重新載入。這就是為什麼有時候切個微信回來，遊戲就沒了。

\booksubsection{儲存（ROM/快閃記憶體）：長期倉庫}

儲存就像你的書櫃，裝所有照片、影片、App的地方。現在主流是UFS 3.1或UFS 4.0，速度比老式的eMMC快好幾倍。

\textbf{為什麼儲存速度重要？}它決定了App開啟速度、遊戲載入速度、照片儲存速度。用久了之後，慢速儲存會讓手機越來越卡。

\booksubsection{螢幕：你的眼睛}

螢幕是你和手機互動的主要視窗，好壞直接影響體驗。

\begin{itemize}
  \item \textbf{材質}：OLED色彩豔麗、對比度高、更省電（顯示黑色時畫素不發光），但低亮度下可能有頻閃傷眼；LCD更護眼、成本低，但對比度和色彩不如OLED。
  \item \textbf{解析度}：1080P是主流，2K更清晰但更耗電。蘋果的"視網膜屏"概念就是：在一定距離下，人眼分辨不出畫素點。
  \item \textbf{重新整理率}：60Hz是標準，90Hz/120Hz更流暢，滑動時絲般順滑。高重新整理率會耗電，所以很多手機支援自適應重新整理率（LTPO），根據內容動態調整。
  \item \textbf{觸控取樣率}：影響觸控響應速度，遊戲玩家要關注，越高越跟手。
  \item \textbf{調光方式}：DC調光比PWM調光更護眼，低亮度下不容易眼睛累。現在很多手機支援高頻PWM調光，也是為了護眼。
\end{itemize}

\booksubsection{螢幕HDR技術}

HDR（High Dynamic Range）讓螢幕同時展現更亮的亮部和更暗的暗部，畫面層次感大幅提升。

\begin{itemize}
  \item \textbf{實現原理}：普通螢幕的亮度範圍有限（SDR），最亮和最暗之間的差距不夠大。HDR螢幕的峰值亮度更高（通常1000尼特以上），對比度更大，色域更廣，能顯示更多細節。
  \item \textbf{主流標準}：HDR10（最基礎，免費開源）、HDR10+（三星推動，動態後設資料）、Dolby Vision（杜比視界，畫質最佳但需授權）。手機支援哪個標準取決於螢幕硬體和系統授權。
  \item \textbf{實際效果}：看HDR影片時，陽光更刺眼、陰影更深邃、色彩更豐富。前提是影片本身是HDR製作的，螢幕也支援HDR顯示。
\end{itemize}

\booksubsection{影像HDR技術}

拍照時的HDR和螢幕HDR是兩回事。影像HDR是透過演算法解決“亮處過曝、暗處死黑”的問題。

\begin{itemize}
  \item \textbf{演算法原理}：手機在按下快門時，實際拍攝了多張不同曝光的照片（欠曝保留亮部細節、過曝保留暗部細節、正常曝光保留中間調），然後將它們合成為一張高動態範圍的照片。
  \item \textbf{處理流程}：多幀捕獲 對齊補償（防手抖） 曝光融合/色調對映 色彩校正 輸出最終照片。
  \item \textbf{效果最佳化}：現在手機已經能做到"實時HDR"——取景器裡看到的畫面就是HDR處理後的效果，所見即所得。蘋果的Smart HDR和安卓的自動HDR都非常成熟。
\end{itemize}

\booksubsection{電池與充電：能量之源}

電池是手機最“短命”的部件，用個兩三年就會明顯老化。現在主流容量是5500-7500mAh。

\begin{itemize}
  \item \textbf{充電功率}：安卓陣營已經卷到了100W甚至200W，十幾分鍾充滿；蘋果還在20-30W徘徊。快充傷電池嗎？用原裝配件問題不大，手機有溫控保護。
  \item \textbf{無線充電}：方便但效率低、發熱大。適合辦公室隨手充，不適合急著出門時用。
  \item \textbf{反向充電}：手機變成充電寶，給耳機、手錶充電，應急用。
\end{itemize}

\booksubsection{攝像頭：光與影的魔法}

現在手機攝像頭動輒三四個，各有分工：

\begin{itemize}
  \item \textbf{主攝}：最常用的，感測器最大，畫質最好。
  \item \textbf{超廣角}：拍風景、建築，視野更寬。
  \item \textbf{長焦}：拍遠處，有光學變焦和數碼變焦之分。
  \item \textbf{微距}：拍極近距離的小東西。
\end{itemize}

\textbf{畫素數不是全部！}感測器尺寸（底大一級壓死人）、光圈大小（f/1.8比f/2.4大）、防抖（OIS光學防抖很重要）都影響畫質。計算攝影（AI演算法）現在越來越重要，能把多張照片合成一張好照片。

\booksubsection{潛望式長焦鏡頭}

傳統長焦鏡頭需要較長的光學路徑，會讓手機變得很厚。潛望式長焦透過一個90度的稜鏡將光線折射，讓光路在手機內部橫向展開，從而在保持手機輕薄的同時實現5倍甚至10倍光學變焦。

\begin{itemize}
  \item \textbf{結構原理}：光線從鏡頭進入 稜鏡折射90度 沿手機寬度方向穿過鏡片組 到達感測器。相當於把一個豎著的望遠鏡橫著放。
  \item \textbf{光學特性}：光學變焦倍數高（5x-10x），畫質遠超數碼變焦。但感測器尺寸通常比主攝小，暗光表現一般。
  \item \textbf{成像優勢}：拍月亮、拍遠處建築、拍舞臺特寫，是普通手機無法企及的。
\end{itemize}

\booksubsection{OIS光學防抖 vs EIS電子防抖}

防抖是影響拍照和錄影畫質的關鍵技術，兩種方式原理完全不同：

\begin{itemize}
  \item \textbf{OIS（光學防抖）}：透過物理移動鏡頭或感測器來抵消手抖。鏡頭內部有微型陀螺儀檢測抖動方向和幅度，然後驅動馬達反向移動鏡片組。優點是不損失畫質，缺點是成本高、佔用空間。
  \item \textbf{EIS（電子防抖）}：透過軟體演算法裁切和偏移畫面來抵消抖動。手機錄製稍大的畫面範圍，然後根據陀螺儀資料裁切出穩定的區域。優點是不需要額外硬體，缺點是會裁切畫面（視野變小），暗光下可能模糊。
  \item \textbf{實際選擇}：拍照優先看OIS，錄影兩者結合效果最好。現在很多旗艦手機支援"Sensor-Shift"感測器位移防抖，比傳統鏡頭防抖效果更好。
\end{itemize}

\booksubsection{光學變焦 vs 數碼變焦}

\begin{itemize}
  \item \textbf{光學變焦}：透過物理移動鏡片組改變焦距，畫質不損失。手機上的光學變焦通常只有固定倍數（如2x、3x、5x、10x），不能像相機那樣連續變焦。
  \item \textbf{數碼變焦}：放大畫面中心區域，本質是裁切+插值放大。放大越多，畫質越差。100倍數碼變焦看著很酷，但實際畫質已經很糊了。
  \item \textbf{混合變焦}：結合光學和數碼，在光學變焦倍數之間用多幀合成演算法彌補畫質損失。效果比純數碼好，但仍不如光學。
\end{itemize}

\booksubsection{音訊：揚聲器與獨顯晶片}

\textbf{單揚聲器 vs 雙揚聲器}：

\begin{itemize}
  \item \textbf{單揚聲器}：只有底部一個發聲單元，聲音偏向一側，沒有立體感。中低端手機常見。
  \item \textbf{雙揚聲器（立體聲）}：頂部和底部各一個揚聲器，形成立體聲場。看影片、聽音樂時左右聲道分明，沉浸感強很多。聽筒兼做頂部揚聲器是常見方案。
  \item \textbf{音質差異}：雙揚聲器不只是"聲音更大"，關鍵是空間感和層次感。用單揚聲器看大片，總覺得少了點什麼。
\end{itemize}

\textbf{遊戲手機獨顯晶片}：

部分遊戲手機和效能手機配備了獨立的顯示處理晶片，它的作用是：

\begin{itemize}
  \item \textbf{插幀}：將30fps的遊戲畫面透過演算法插值到60fps或120fps，畫面更流暢。這個過程在獨顯晶片上完成，不佔用SoC的GPU資源。
  \item \textbf{超分}：將低解析度畫面透過AI演算法提升到高解析度，畫質更清晰。
  \item \textbf{降功耗}：因為插幀和超分由獨顯晶片負責，SoC的負載降低，發熱和功耗都減少，遊戲更穩定。
\end{itemize}

\booksubsection{GPS定位：單頻與雙頻}

手機定位主要依賴GPS等衛星導航系統。

\begin{itemize}
  \item \textbf{單頻GPS}：只接收L1頻段訊號，定位精度約3-5米。在開闊地帶夠用，但在高樓密集的城市峽谷中，訊號反射會導致定位偏移較大。
  \item \textbf{雙頻GPS}：同時接收L1和L5兩個頻段訊號。L5頻段波長更長、抗干擾能力更強，雙頻聯合解算可以消除大部分電離層誤差，定位精度可達1-3米。在城市中導航、打車定位時優勢明顯。
  \item \textbf{多星座支援}：現代手機不僅支援GPS，還支援北斗（中國）、GLONASS（俄羅斯）、Galileo（歐洲）等多個衛星系統。同時連線的衛星越多，定位越快越準。
\end{itemize}

\booksubsection{感測器：手機的小觸角}

手機裡藏著十幾個感測器，默默工作：

\begin{itemize}
  \item \textbf{加速度計+陀螺儀}：感知手機姿態，玩遊戲轉向、自動旋轉螢幕。
  \item \textbf{光線感測器}：自動調節螢幕亮度。
  \item \textbf{距離感測器}：打電話時手機貼近耳朵自動熄屏。
  \item \textbf{指紋感測器}：現在主流是屏下光學或超聲波指紋。
  \item \textbf{面部識別}：蘋果的Face ID是3D結構光，安全級別高；很多安卓是2D照片識別，安全性一般。
  \item \textbf{NFC}：近場通訊，刷公交卡、門禁卡、銀行卡。
  \item \textbf{氣壓計}：測海拔高度，爬樓梯都能檢測到。
  \item \textbf{霍爾感測器}：檢測磁吸殼的開合。
\end{itemize}

\booksection{二、作業系統：安卓 VS iOS VS 鴻蒙}

硬體是身體，作業系統是靈魂。三大系統各有特色：

\begin{itemize}
  \item \textbf{安卓}：像開源社群，手機廠商可以深度定製，選擇多。你可以換主題、裝第三方應用商店、折騰各種玩法。但也因為開放，碎片化嚴重，不同品牌體驗差異大。
  \item \textbf{iOS}：像精心打理的花園，蘋果軟硬體一手抓，體驗統一流暢。應用質量高，隱私保護好，但自由度低，很多事得按蘋果的規矩來。
  \item \textbf{鴻蒙}：華為的新物種，主打分散式能力，可以讓不同裝置無縫協同。後面我們會專門講。
\end{itemize}

三者生態都極其豐富，選你喜歡的生活方式就好。沒有絕對的好壞，只有適不適合。

\booksection{三、手機介面全解析}

手機介面是連線外部世界的橋樑，不同介面差異很大：

\booksubsection{Micro USB}

\begin{itemize}
  \item \textbf{物理特性}：扁平梯形介面，分正反面（插反了插不進去）。已經在手機上基本淘汰，但部分低端裝置和老裝置仍在使用。
  \item \textbf{傳輸限制}：僅支援USB 2.0標準，最高傳輸速率480Mbps，充電功率通常不超過10W。
  \item \textbf{缺點}：正反面不通用、插拔壽命短、傳輸速度慢、充電功率低。
\end{itemize}

\booksubsection{Lightning（蘋果閃電介面，已淘汰）}

\begin{itemize}
  \item \textbf{物理特性}：蘋果自研的8針介面，正反都可以插。iPhone 14及更早機型使用，現已全面淘汰。
  \item \textbf{傳輸限制}：Lightning介面物理上僅支援USB 2.0標準，最高傳輸速率480Mbps。即使使用USB-C轉Lightning線，傳輸速度也受限於Lightning介面本身。
  \item \textbf{淘汰原因}：歐盟透過法規要求2024年起所有手機統一使用USB-C介面，蘋果從iPhone 15系列開始全面改用Type-C。Lightning介面正式退出歷史舞臺。
  \item \textbf{注意}：蘋果MFi認證的資料線內含認證晶片，非認證線可能無法正常充電或傳輸資料。如果你仍在使用舊款iPhone，建議備好Lightning資料線。
\end{itemize}

\booksubsection{USB Type-C}

\begin{itemize}
  \item \textbf{物理特性}：24針腳的橢圓形介面，正反都可以插。已成為手機和多數電子裝置的標準介面。歐盟已立法要求2024年起所有手機、平板、耳機等便攜裝置統一使用USB-C介面，全球其他地區也在跟進。
  \item \textbf{傳輸協議}：Type-C只是介面形態，實際傳輸能力取決於內部協議。支援USB 2.0（480Mbps）、USB 3.0/3.1/3.2（5-20Gbps）甚至USB4（40Gbps）。
  \item \textbf{重要區分}：並非所有Type-C介面都支援USB 3.0！很多手機雖然用Type-C介面，但內部只連線了USB 2.0的引腳，傳輸速度和老介面一樣慢。iPhone 16/17標準版僅支援USB 2.0速度，Pro版才支援USB 3.0（需使用USB-C to USB-C線）。
  \item \textbf{充電能力}：支援USB PD等快充協議，最高可支援240W充電。
  \item \textbf{影片輸出}：支援DisplayPort Alt Mode，可以透過轉接線輸出影片到顯示器。
\end{itemize}

\booksubsection{5mm音訊介面}

\begin{itemize}
  \item \textbf{物理特性}：圓形3段或4段插孔，3段傳輸左右聲道音訊，4段多一個麥克風通道。
  \item \textbf{特點}：模擬訊號傳輸，即插即用，不需要驅動或轉換。音質穩定，延遲為零。
  \item \textbf{現狀}：大部分旗艦手機已取消3.5mm介面，需要Type-C轉接線或使用藍芽耳機。部分中端機和遊戲手機仍保留。
\end{itemize}

\booksection{四、手機產品體系與命名規則}

面對各品牌眼花繚亂的產品線，搞清楚命名規則能幫你快速定位適合自己的機型。

\booksubsection{主流品牌產品系列：}

\begin{itemize}
  \item \textbf{蘋果}：iPhone（標準版）/ iPhone Plus（大屏版）/ iPhone Pro（專業版）/ iPhone Pro Max（大屏專業版）/ iPhone SE（平價版）/ iPhone 16e（平價AI版，搭載A18晶片和蘋果自研C1基帶）。2026年iPhone 17系列全面搭載A19 Pro晶片，AI能力進一步提升。
  \item \textbf{華為}：Mate系列（商務旗艦）/ P/Pura系列（影像旗艦）/ Mate X系列（摺疊旗艦）/ nova系列（年輕中端）/ 暢享系列（入門）
  \item \textbf{小米}：小米數字系列（旗艦）/ 小米MIX（探索旗艦/摺疊）/ Redmi K系列（效能旗艦）/ Redmi Note系列（中端）/ Redmi數字系列（入門）
  \item \textbf{OPPO}：Find X系列（旗艦）/ Reno系列（中端）/ A系列（入門）/ 一加數字系列（效能旗艦）
  \item \textbf{vivo}：X系列（影像旗艦）/ iQOO數字系列（效能旗艦）/ S系列（自拍中端）/ Y系列（入門）
  \item \textbf{三星}：Galaxy S系列（旗艦）/ Galaxy Z Fold/Flip（摺疊旗艦）/ Galaxy A系列（中端）
\end{itemize}

\booksubsection{常見字尾含義：}

\begin{itemize}
  \item \textbf{Pro}：Professional的縮寫，專業版。通常在標準版基礎上升級攝像頭、螢幕、處理器等核心配置。
  \item \textbf{Max}：最大版。螢幕最大、電池最大，但配置不一定比Pro高。
  \item \textbf{Ultra}：極致版。通常代表該系列的最高配置，堆料最猛，價格也最高。如三星S24 Ultra、小米14 Ultra。
  \item \textbf{SE}：Special Edition特別版，實際多為平價版。配置縮水，價格降低。如iPhone SE。
  \item \textbf{Plus}：加大版。螢幕比標準版大，電池也大，但配置基本一致。
  \item \textbf{Lite/青春版}：輕量版。配置精簡，價格更低。
  \item \textbf{Note}：原本指大屏手機，現在各品牌用法不一。Redmi Note是中端系列。
  \item \textbf{Neo}：新銳版。通常指效能取向的中端機型。
\end{itemize}

\booksection{五、基帶與天線：手機通訊的幕後功臣}

手機能打電話、上網，靠的是基帶晶片和天線這對搭檔。

\booksubsection{基帶晶片：}

基帶是手機中負責無線通訊的處理器，它將數字訊號調製成無線電波發出去，也將接收到的無線電波解調成數字訊號。基帶效能直接影響通話質量、上網速度和訊號穩定性。高通是最大的手機基帶供應商，聯發科、三星也有自研基帶。蘋果從iPhone 16e開始搭載自研的C1基帶晶片，標誌著蘋果正式進入自研基帶時代，C1基帶採用4nm製程，功耗表現優異；iPhone 16/16 Plus/16 Pro/16 Pro Max仍使用高通基帶（驍龍X75 5G），預計後續機型將逐步過渡到自研基帶。更早的iPhone XS/11系列使用Intel基帶時訊號表現不佳，是不少使用者的痛點。

\booksubsection{天線設計：}

手機天線藏在機身內部，通常圍繞邊框佈置。天線設計直接影響訊號接收效能：

\begin{itemize}
  \item \textbf{MIMO技術}：多入多出，使用多個天線同時收發資料，提升速度和穩定性。4×4 MIMO比2×2 MIMO速度更快。
  \item \textbf{天線數量}：旗艦機通常有十幾個天線（蜂窩、Wi-Fi、藍芽、NFC、GPS各需要獨立天線）。
  \item \textbf{金屬機身的影響}：金屬會遮蔽訊號，所以金屬機身手機需要設計天線帶（邊框上的塑膠條）。玻璃和塑膠機身對訊號影響較小。
\end{itemize}

\booksubsection{網路技術演進：}

\begin{itemize}
  \item \textbf{4G LTE}：理論峰值速率下行1Gbps/上行500Mbps（Cat.18），實際體驗通常50-150Mbps。覆蓋最廣，夠用。
  \item \textbf{5G NR}：理論峰值速率下行10Gbps+。Sub-6GHz頻段覆蓋較好，mmWave毫米波速度極快但覆蓋極小。
  \item \textbf{5G-A（5G Advanced）}：5G的增強版，也叫5.5G。下行速率可達10Gbps，上行1Gbps，連線數密度和定位精度大幅提升。2024年起中國移動、電信、聯通開始在重點城市部署5G-A網路，華為Mate 70/80系列、小米15/16系列等旗艦手機已支援5G-A。5G-A還帶來了"通感一體"能力，基站可以同時通訊和感知環境，為低空經濟（無人機配送）和車聯網提供基礎設施。
\end{itemize}

\booksubsection{網路頻段分類：}

手機支援的網路頻段決定了它能在哪些國家、哪些運營商的網路上使用：

\begin{itemize}
  \item \textbf{FDD-LTE Band（B系列）}：如B1、B3、B5等，4G頻段。不同國家使用不同頻段，中國主要用B1/B3/B5/B8/B34/B38/B39/B40/B41。
  \item \textbf{TDD-LTE Band}：4G的另一制式，中國大量使用B38/B39/B40/B41。
  \item \textbf{5G NR Band（N系列）}：如n1、n3、n28、n41、n78、n79等。n78（3.5GHz）是中國5G主力頻段，n41（2.5GHz）移動使用較多。
  \item \textbf{購買建議}：如果需要出國使用，購買前確認手機是否支援目標國家的頻段。國行手機通常對中國頻段支援最全，但可能缺少部分海外頻段。
\end{itemize}

\booksection{六、你的手機，到底在扮演多少角色？}

它早已不只是電話：

\begin{enumerate}
  \item \textbf{通訊工具}：電話、簡訊、微信、郵件……
  \item \textbf{智慧相機}：拍照、錄影、掃描文件……
  \item \textbf{錢包}：支付、公交卡、銀行卡、會員卡……
  \item \textbf{資訊門戶}：新聞、天氣、股票、知識……
  \item \textbf{遊戲機}：從消消樂到3A大作……
  \item \textbf{移動辦公室}：文件、表格、會議、協作……
  \item \textbf{健康助手}：計步、心率、睡眠、運動……
  \item \textbf{智慧家居遙控器}：控制燈、空調、電視、掃地機器人……
  \item \textbf{數字身份證}：電子身份證、電子駕照、健康碼……
  \item \textbf{娛樂中心}：音樂、影片、電子書、播客……
\end{enumerate}

想想看，一臺小小的裝置，承載了多少生活！所以，學會用好它，就是學會更好地生活。

\begin{tcbpractice}{實操一下：瞭解硬體配置}
在 [設定] 裡，看看能不能找到處理器型號、執行記憶體（RAM）和機身儲存（ROM）的具體大小。對你手機的能力有個數。安卓使用者可以下載“安兔兔”或“DevCheck”檢視詳細硬體資訊。

\end{tcbpractice}
